% /Users/pikachu/books/payments-book/src/parts/part6/ch31-source-map.tex
\chapter{Карта первоисточников}
\label{ch:primary-sources}
\margintoc


\begin{kaobox}[title=Что вы узнаете из этой главы]
  Как быстро понять, какой документ является owner-источником
  по конкретной теме; чем отличаются public и restricted документы;
  в каком порядке входить в EMVCo, PCI SSC, правила карточных сетей,
  документы НСПК, Банка России и ISO; и как не тратить часы
  на чтение вторичных пересказов вместо нормативного текста.
\end{kaobox}

Когда платёжная команда говорит \enquote{надо проверить в правилах},
обычно неясно, о каких именно правилах идёт речь. У одной темы могут
существовать сразу три слоя документов: публичный бизнес-правила,
техническая спецификация и закрытые operational bulletins. Поэтому
первая задача навигации — не \enquote{найти документ вообще},
а понять, какой документ является owner-источником для данного вопроса.

\begin{table}[h]
\centering
\scriptsize
\renewcommand{\arraystretch}{1.15}
\begin{tabular}{|p{1.8cm}|p{3.3cm}|p{2.0cm}|p{3.0cm}|}
\hline
\textbf{Организация} & \textbf{Ключевые документы} & \textbf{Доступ} & \textbf{Когда читать} \\
\hline
EMVCo & Books 1--4, Contactless, Tokenisation, 3DS & mostly public & когда нужен протокол карты, терминала или токена \\
\hline
PCI SSC & PCI DSS, SAQ, PIN Security, P2PE & public & когда вопрос про scope, контроль или аудит \\
\hline
Visa & Core Rules, VCR guides, TADG, BASE I/II, VTS docs & mixed & когда вопрос про обязанности участников или VisaNet flow \\
\hline
Mastercard & Mastercard Rules, Chargeback Guide, Banknet/IPM, MDES & mixed & когда нужен Banknet, IPM или policy Mastercard \\
\hline
НСПК / Мир & Правила Мир, SBP docs, техспеки НСПК & mostly restricted & когда вопрос про domestic rail, Мир или СБП \\
\hline
Банк России & 161-ФЗ, 762-П, 382-П, 683-П, 266-П & public & когда вопрос про правовую рамку и обязательные требования \\
\hline
ISO & ISO 8583, ISO 20022, ISO 4217, ISO 9564 & mixed / paid & когда нужен язык стандарта, а не только сетевой диалект \\
\hline
\end{tabular}
\end{table}

Эта таблица не заменяет чтение документов, но помогает выбрать
правильную точку входа. Ниже — не каталог ссылок, а метод
\enquote{что читать первым и зачем}.

\section{EMVCo}
\label{sec:ps-emvco}

EMVCo — owner-источник по темам, где важен протокол взаимодействия
карты, терминала, токена и 3DS-доменов. Если вопрос звучит как
\enquote{что именно обязан сделать терминал / карта / wallet},
искать нужно здесь, а не в правилах эквайера.

\subsection*{Какие документы являются основными}

\textbf{EMV Chip Specifications} — Books 1--4; именно они задают
модель контактной EMV-транзакции (см. главу \ref{ch:emv}).
\textbf{EMV Contactless Specifications} — owner-источник для NFC
и kernel-логики (см. главу \ref{ch:contactless}).
\textbf{EMV Payment Tokenisation Specification} — owner-документ
по ролям TSP, token domain controls и token lifecycle
(см. главу \ref{ch:tokenization}).
\textbf{EMV 3-D Secure Protocol and Core Functions Specification} —
owner-источник по AReq/ARes, challenge flow и message-level
логике 3DS (см. главу \ref{ch:3ds}).

\subsection*{Как входить}

Для первого чтения полезно начинать не с Book 1, а с того раздела,
который отвечает на конкретный вопрос. Если нужен жизненный цикл
карточной транзакции — Book 3. Если нужен UX терминала и кассира —
Book 4. Если нужен contactless entry point — Book A/B/C.
EMVCo плохо подходит для линейного чтения от корки до корки,
но отлично работает как owner-документ для точечной проверки.

\section{PCI SSC}
\label{sec:ps-pci}

PCI SSC — owner-источник по вопросам scope и контролей.
Если команда спорит, можно ли использовать SAQ A, обязан ли мерчант
контролировать JavaScript на checkout-странице или как защищается
PIN block, спор нужно закрывать документами PCI, а не блог-постами PSP.

\subsection*{Какие документы являются основными}

\textbf{PCI DSS} — главный документ по контролям среды,
где появляются данные держателя карты (см. главу \ref{ch:pci-dss}).
\textbf{SAQ} — сокращённые owner-документы для мерчантов
по типам интеграции. \textbf{PCI PIN Security Requirements} —
owner-источник по PIN и HSM-контру. \textbf{PCI P2PE} —
owner-документ, если вопрос про шифрование от терминала
до точки дешифрования.

\subsection*{Как входить}

Для практической работы почти всегда полезно идти сверху вниз:
сначала определить scope, потом выбрать применимый SAQ/ROC-путь,
и только после этого читать конкретные требования. Ошибка
в определении scope почти всегда дороже, чем ошибка в трактовке
одного подпункта.

\section{\scheme{Visa}}
\label{sec:ps-visa}

Документация \scheme{Visa} — это два параллельных мира:
публичные правила о том, что разрешено участникам,
и ограниченные технические документы о том, как именно выглядит
реализация в VisaNet.

\subsection*{Какие документы являются основными}

\textbf{Visa Core Rules and Visa Product and Service Rules} —
owner-источник по обязанностям эмитента, эквайера и мерчанта.
\textbf{Visa Dispute Management Guidelines} — owner-документ
по reason codes, representment и dispute-policy.
\textbf{Visa Technical Acceptance Device Guide (TADG)} —
технический owner-документ по приёму карт на терминалах.
\textbf{BASE I / BASE II} — owner-источник по авторизации
и клирингу на уровне VisaNet. \textbf{VTS documentation} —
owner-документы по Visa Token Service.

\subsection*{В каком порядке читать}

Если вопрос policy-level — сначала Core Rules. Если question-level
спор касается технического поля, терминала или acceptance behavior —
ищите TADG или техническую спецификацию через участника сети.
Developer-портал \domain{developer.visa.com} полезен для API-продуктов,
но не заменяет сетевые правила.

\section{\scheme{Mastercard}}
\label{sec:ps-mastercard}

\scheme{Mastercard} организует документы похожим образом:
сначала правила и chargeback-политика, затем ограниченные
Banknet/IPM-спецификации и developer-материалы.

\subsection*{Какие документы являются основными}

\textbf{Mastercard Rules} — owner-источник по policy-level
обязательствам. \textbf{Mastercard Chargeback Guide} —
owner-документ по dispute flow. \textbf{Banknet Processing
and IPM Technical Manual} — owner-источник по авторизации
и клирингу через Banknet. \textbf{MDES} — owner-документация
по токенизации. \textbf{M/Chip Requirements} — owner-документ
по Mastercard-специфике EMV.

\subsection*{В каком порядке читать}

Начинать полезно с Rules, если вопрос касается обязанностей
участников, и с IPM / Banknet manual, если нужен протокол
или клиринговая модель. Developer-портал Mastercard помогает
с открытыми API, но не заменяет restricted docs.

\section{НСПК и Мир}
\label{sec:ps-nspk}

НСПК — owner-источник по domestic card rail \scheme{Мир},
внутрироссийскому эквайрингу и СБП. Важная особенность:
бизнес-логика частично публична, а большая часть технической
документации остаётся ограниченной.

\subsection*{Какие документы являются основными}

\textbf{Правила платёжной системы Мир} — owner-источник
по ролям и обязанностям участников. \textbf{Спецификация
бесконтактных платежей Мир} — owner-документ по Mir Contactless.
\textbf{Технические спецификации НСПК} — owner-документы
по форматам сообщений, подключению и МРА. \textbf{Документация
по СБП} — owner-источник по API и operational-правилам СБП.

\subsection*{Как входить}

Если вопрос правовой или продуктовый, обычно хватает публичных
правил. Если вопрос касается message format, сертификации
или тонкой логики domestic processing, почти наверняка понадобится
restricted documentation через участника системы.

\section{Банк России}
\label{sec:ps-cbr}

Банк России — owner-источник для правовой рамки, а не для message-level
протокола. Когда вопрос звучит как \enquote{обязаны ли мы},
\enquote{имеем ли мы право} или \enquote{какой статус нужен},
проверять надо нормативные акты ЦБ и федеральные законы.

\subsection*{Какие документы являются основными}

\law{161-ФЗ} задаёт базовые роли платёжного рынка.
\law{762-П} описывает правила перевода денежных средств.
\law{382-П} и \law{683-П} задают требования по защите информации.
\law{266-П} добавляет карточную специфику. Для цифрового рубля,
СБП и иных специальных тем owner-документы будут свои, но логика
поиска та же: начинать с нормативного текста, а не с пересказа.

\subsection*{Как входить}

Читать нормативные акты полезно вместе: закон задаёт роли,
положение — operational-логику, письма и разъяснения показывают,
как регулятор трактует спорные места. Для практической работы
удобнее использовать правовые системы с историей изменений
и перекрёстными ссылками.

\section{ISO}
\label{sec:ps-iso}

ISO — это слой, где платёжная индустрия разговаривает на языке
стандартов, но почти никогда не работает по \enquote{чистому ISO}
без сетевого диалекта.

\subsection*{Какие документы являются основными}

\iso{8583} — owner-стандарт по структуре financial messages;
на практике его всегда дополняют сетевые реализации.
\iso{20022} — owner-стандарт по современной модели финансовых
сообщений, критичный для cross-border, SWIFT MX и многих
instant payment rail (см. главу \ref{ch:multicurrency}).
ISO 4217, ISO 3166 и ISO 9564 — owner-документы по кодам валют,
кодам стран и PIN block.

\subsection*{Как входить}

Если задача прикладная и привязана к конкретной сети, сначала
читайте сетевой диалект, затем ISO как верхний слой терминов
и модели. Если задача архитектурная или межсистемная,
\iso{20022} стоит читать сразу как owner-модель.

\begin{chaptersummary}
\begin{itemize}
  \item У каждой темы должен быть owner-документ: EMVCo для
        card / token / 3DS protocol, PCI SSC для controls,
        сети для obligations и their rails, ЦБ для правовой рамки,
        ISO для языка стандартов.
  \item Разделяйте public и restricted документы заранее:
        это экономит время и не создаёт ложного ожидания,
        что любой вопрос можно закрыть публичной PDF.
  \item Начинать чтение стоит с документа, который отвечает
        на ваш тип вопроса: policy, protocol, compliance
        или regulation.
  \item Developer-порталы полезны для API-интеграции, но редко
        являются owner-источником по сетевым правилам.
  \item Хорошая навигация начинается с карты источников:
        кто владеет истиной по теме, где документ доступен
        и в какой последовательности его читать.
\end{itemize}
\end{chaptersummary}

